\documentclass{../nndlcourse}

\begin{document}

\title{线性回归模型}
\author{数据科学与大数据技术专业}
\date{第 2 讲 }
\maketitle



\section{概率论基础回顾}

\begin{enumerate}
  \item 概率空间与随机变量
    \begin{itemize}
      \item 概率空间记为 $(\Omega, \mathcal{F}, P)$, 用于刻画随机试验及其不确定性.
        \begin{itemize}
          \item 样本空间 $\Omega$: 随机试验所有可能结果的集合.
          \item 样本点 $\omega$: $\Omega$ 中的一个基本结果.
          \item 事件 $A$: 样本空间的子集.
          \item 事件域 $\mathcal{F}$: 定义在 $\Omega$ 上的 $\sigma$-代数.
          \item 概率测度 $P$: 满足非负性、规范性和可列可加性的函数.
        \end{itemize}
      \item 随机变量是定义在概率空间上的可测函数, 常用来将随机结果映射为数值.
      \item 随机变量的分布通常通过以下对象描述:
        \begin{itemize}
          \item 分布函数(CDF): $F(x) = P(X \le x)$.
          \item 概率质量函数(PMF): 用于离散型随机变量, 给出 $P(X = x)$.
          \item 概率密度函数(PDF): 用于连续型随机变量, 满足 $P(a \le X \le b) = \int_a^b f(x) \, dx$.
        \end{itemize}
      \item 按取值形式, 随机变量可分为离散型和连续型; 前者常见于计数问题, 后者常见于测量问题.
    \end{itemize}
  \item 随机向量与分布关系
    \begin{itemize}
      \item 随机向量 $\boldsymbol{X} = (X_1, X_2, \dots, X_n)$ 用于同时描述多个随机变量.
      \item 联合分布刻画多个随机变量同时取值的概率规律; 对连续型随机向量, 可进一步用联合密度函数描述.
      \item 边际分布可由联合分布对其余变量求和或积分得到, 因而只反映单个变量的统计规律.
      \item 条件概率刻画在已知部分信息时其余事件发生的可能性, 贝叶斯公式为
        \begin{equation*}
        P(A \mid B) = \frac{P(B \mid A) P(A)}{P(B)}.
        \end{equation*}
      \item 若 $X$ 与 $Y$ 的联合分布等于各自边际分布的乘积, 则称它们相互独立.
    \end{itemize}
  \item 采样、期望与大数定律
    \begin{itemize}
      \item 常见采样方法包括:
        \begin{itemize}
          \item 直接采样: 从已知分布直接生成样本.
          \item 逆变换采样: 先生成均匀随机变量, 再利用分布函数的反函数得到目标分布样本.
        \end{itemize}
      \item 期望反映随机变量在重复试验下的平均水平, 是许多学习准则和统计量的基础.
      \item 大数定律说明: 当样本量充分大时, 样本均值会以概率收敛的方式逼近期望, 这为用样本统计量近似总体统计量提供了理论依据.
    \end{itemize}
\end{enumerate}

\textbf{问题思考}: 

\begin{itemize}
  \item 什么是概率空间? 请举例说明. 
  \item 已知一个二维随机向量的边缘分布, 能否知道它的联合分布? 
  \item 如何理解大数定律在实际中的意义? 
\end{itemize}


% ===================== 线性回归模型 =====================
\section{线性回归概述}
\begin{enumerate}
  \item \textbf{模型表示}
    \begin{itemize}
    \item 线性回归假设输出是输入特征的线性函数. 若将偏置项并入增广特征向量, 则模型可写为
    \begin{equation*}
    \hat{y} = f(\boldsymbol{x}; \boldsymbol{w}) = \boldsymbol{w}^{\top} \boldsymbol{x}.
    \end{equation*}
    \item 其中 $\boldsymbol{x} \in \mathbb{R}^M$ 表示样本特征, $\boldsymbol{w} \in \mathbb{R}^M$ 表示待学习参数.
    \item 该模型主要用于连续值预测任务, 如房价预测、销量预测和温度估计等.
    \end{itemize}
  \item \textbf{基本特点}
  \begin{itemize}
    \item 模型形式简单、参数含义清晰, 便于分析和解释.
    \item 在线性假设成立或经过特征变换后近似成立时, 线性回归往往能够给出有效基线.
  \end{itemize}
\end{enumerate}

\section{参数估计方法}

\begin{enumerate}
    \item \textbf{经验风险最小化}
    \begin{itemize}
    \item 在线性回归中, 常用平方损失作为学习准则. 对于训练集 $\{(\boldsymbol{x}^{(n)}, y^{(n)})\}_{n=1}^N$, 经验风险可写为
        \begin{align*}
    \mathcal{R}(\boldsymbol{w}) & = \sum_{n=1}^N \mathcal{L}\left(y^{(n)}, f\left(\boldsymbol{x}^{(n)}; \boldsymbol{w}\right)\right) \\
    & = \frac{1}{2} \sum_{n=1}^N \left(y^{(n)} - \boldsymbol{w}^{\top} \boldsymbol{x}^{(n)}\right)^2.
        \end{align*}

    \item 若将全部样本按列堆叠成矩阵
    \begin{equation*}
      \boldsymbol{X} = \left[\boldsymbol{x}^{(1)}, \boldsymbol{x}^{(2)}, \cdots, \boldsymbol{x}^{(N)}\right] \in \mathbb{R}^{M \times N},
      \quad
      \boldsymbol{y} = \left[y^{(1)}, y^{(2)}, \cdots, y^{(N)}\right]^{\top} \in \mathbb{R}^{N},
    \end{equation*}
    则预测向量满足
    \begin{equation*}
      \hat{\boldsymbol{y}} = \boldsymbol{X}^{\top} \boldsymbol{w}.
    \end{equation*}

    \item 因而经验风险可进一步写成紧凑的矩阵形式
    \begin{equation*}
    \mathcal{R}(\boldsymbol{w}) = \frac{1}{2} \left\| \boldsymbol{y} - \boldsymbol{X}^{\top} \boldsymbol{w} \right\|^2.
    \end{equation*}

    \item \textbf{解析解的推导要点}:
        \begin{itemize}
      \item 对目标函数求梯度可得
            \begin{equation*}
        \frac{\partial \mathcal{R}(\boldsymbol{w})}{\partial \boldsymbol{w}} = -\boldsymbol{X}\left(\boldsymbol{y} - \boldsymbol{X}^{\top} \boldsymbol{w}\right) = \boldsymbol{X}\boldsymbol{X}^{\top} \boldsymbol{w} - \boldsymbol{X}\boldsymbol{y}.
            \end{equation*}
      \item 令梯度为零, 可得正规方程
            \begin{equation*}
        \boldsymbol{X}\boldsymbol{X}^{\top} \boldsymbol{w} = \boldsymbol{X}\boldsymbol{y}.
            \end{equation*}
      \item 当 $\boldsymbol{X}\boldsymbol{X}^{\top}$ 可逆时, 经验风险最小化的解析解为
            \begin{equation*}
        \boldsymbol{w}^* = \left(\boldsymbol{X}\boldsymbol{X}^{\top}\right)^{-1} \boldsymbol{X}\boldsymbol{y}.
            \end{equation*}
      \item 上式说明: 在线性回归中, 最小二乘解本质上就是平方损失下的经验风险最小化解.
        \end{itemize}
    \end{itemize}

    \item \textbf{结构风险最小化}
    \begin{itemize}
    \item 当特征存在共线性、样本不足或噪声较大时, 矩阵 $\boldsymbol{X}\boldsymbol{X}^{\top}$ 可能不可逆或病态, 从而使经验风险最小化的解析解不稳定.
    \item 一个自然的改进方式是在经验风险中加入参数范数惩罚项, 即采用结构风险最小化:
        \begin{equation*}
        \mathcal{R}_{srm}(\boldsymbol{w}) = \frac{1}{2}\left\|\boldsymbol{y}-\boldsymbol{X}^{\top} \boldsymbol{w}\right\|^2 + \frac{\lambda}{2}\|\boldsymbol{w}\|^2
        \end{equation*}
    其中 $\lambda > 0$ 为正则化系数.

    \item \textbf{加入 $\lambda \|\boldsymbol{w}\|^2$ 的主要作用}:
        \begin{itemize}
      \item 在数值计算上, 需要求逆的矩阵变为 $\boldsymbol{X}\boldsymbol{X}^{\top} + \lambda \boldsymbol{I}$, 可显著改善可逆性和稳定性.
      \item 在统计学习上, 它通过限制参数幅度来控制模型复杂度, 从而缓解过拟合.
        \end{itemize}

    \item 相应的最优解为
    \begin{equation*}
    \boldsymbol{w}^* = \left(\boldsymbol{X}\boldsymbol{X}^{\top} + \lambda \boldsymbol{I}\right)^{-1} \boldsymbol{X}\boldsymbol{y}.
    \end{equation*}
    \end{itemize}
\end{enumerate}



\textbf{问题思考}: 
\begin{itemize}
    \item 为什么线性回归模型需要引入正则化? 
    \item 线性回归的概率解释与最小二乘法有何联系? 
\end{itemize}
    

% ===================== 概率视角解释 =====================
\section{概率视角解释}

从概率建模的角度看, 线性回归不仅是在拟合一个函数, 也可以理解为在学习条件分布 $p(y \mid \boldsymbol{x}; \boldsymbol{w})$. 在这一视角下, 最小二乘法、最大似然估计和带正则化的参数估计之间具有统一的解释.

\begin{enumerate}
  \item \textbf{函数建模与条件概率建模}
    \begin{itemize}
    \item 函数建模关注的是寻找映射关系 $f(\boldsymbol{x}; \boldsymbol{w})$, 并通过最小化损失函数来学习参数.
    \item 条件概率建模则假设数据由某个概率过程生成, 目标是估计给定输入时输出的条件分布.
    \item 对线性回归而言, 这两种视角在高斯噪声假设下会导向一致的参数估计结果.
  \end{itemize}

  \item \textbf{高斯噪声假设与似然函数}
  \begin{itemize}
    \item 假设在给定输入 $\boldsymbol{x}$ 和参数 $\boldsymbol{w}$ 时, 输出满足
    \begin{equation*}
    y \mid \boldsymbol{x}, \boldsymbol{w} \sim \mathcal{N}(\boldsymbol{w}^{\top} \boldsymbol{x}, \sigma^2).
    \end{equation*}
    \item 这里的“概率”强调在参数已知时观测数据出现的可能性; “似然”则是把观测数据固定后, 将其看作参数 $\boldsymbol{w}$ 的函数.
    \item 对于独立同分布样本, 似然函数为
        \begin{equation*}
        L(\boldsymbol{w}) = \prod_{n=1}^N \frac{1}{\sqrt{2\pi}\sigma} \exp\left(-\frac{(y^{(n)} - \boldsymbol{w}^\top \boldsymbol{x}^{(n)})^2}{2\sigma^2}\right)
    \end{equation*}
    \end{itemize}

    \item \textbf{最大似然估计(MLE)}
    \begin{itemize}
    \item 最大似然估计的目标是寻找使对数似然最大的参数 $\boldsymbol{w}_{MLE}^*$.
    \item 去掉与参数无关的常数项后, 负对数似然为
        \begin{equation*}
        -\ln L(\boldsymbol{w}) \propto \frac{1}{2\sigma^2} \sum_{n=1}^N \left(y^{(n)} - \boldsymbol{w}^\top \boldsymbol{x}^{(n)}\right)^2 
        \end{equation*}
    这表明: 当噪声服从高斯分布时, 最大化似然等价于最小化平方误差, 因而与经验风险最小化共享同一个最优解.
        \begin{equation*}
        \boldsymbol{w}^*_{MLE} = \arg\max_{\boldsymbol{w}} \ln L(\boldsymbol{w}) = (\boldsymbol{X}\boldsymbol{X}^\top)^{-1}\boldsymbol{X}\boldsymbol{y}
        \end{equation*}
    \end{itemize}

    \item \textbf{贝叶斯估计}
    \begin{itemize}
    \item 贝叶斯方法不再把参数 $\boldsymbol{w}$ 视为固定常数, 而是将其视为随机变量, 进而学习后验分布 $p(\boldsymbol{w} \mid \boldsymbol{X}, \boldsymbol{y})$.
    \item 若先验取为零均值高斯分布
    \begin{equation*}
    \boldsymbol{w} \sim \mathcal{N}(\boldsymbol{0}, \nu^2 \boldsymbol{I}),
    \end{equation*}
    则后验对数在省略常数项后满足
        \begin{align*}
        \log p(\boldsymbol{w} \mid \boldsymbol{X}, \boldsymbol{y} ; \nu, \sigma) & \propto \log p(\boldsymbol{y} \mid \boldsymbol{X}, \boldsymbol{w} ; \sigma)+\log p(\boldsymbol{w} ; \nu) \\
    & \propto -\frac{1}{2 \sigma^2} \sum_{n=1}^N \left(y^{(n)}-\boldsymbol{w}^{\top} \boldsymbol{x}^{(n)}\right)^2 - \frac{1}{2 \nu^2} \boldsymbol{w}^{\top} \boldsymbol{w} \\
    & = -\frac{1}{2 \sigma^2}\left\|\boldsymbol{y}-\boldsymbol{X}^{\top} \boldsymbol{w}\right\|^2 - \frac{1}{2 \nu^2} \boldsymbol{w}^{\top} \boldsymbol{w}.
        \end{align*}
    \item 后验分布综合了数据提供的信息和参数先验的信息, 因而不仅给出参数点估计, 也刻画了参数不确定性.
    \item 对于新样本 $\boldsymbol{x}^*$, 贝叶斯预测需要对参数后验做积分:
        \begin{equation*}
      p(y^* \mid \boldsymbol{x}^*, \boldsymbol{X}, \boldsymbol{y}) = \int p(y^* \mid \boldsymbol{x}^*, \boldsymbol{w}) p(\boldsymbol{w} \mid \boldsymbol{X}, \boldsymbol{y}) d\boldsymbol{w}.
        \end{equation*}
    \end{itemize}

    \item \textbf{最大后验估计(MAP)与正则化}
    \begin{itemize}
    \item 严格的贝叶斯预测需要处理整个后验分布, 计算量往往较大. 在实践中, 常采用后验分布的众数作为点估计, 即最大后验估计.
    \item 最大化后验概率等价于最小化负对数后验, 因而等价于最小化“平方误差项 $+$ 参数范数项”.
    \item 若记
    \begin{equation*}
    \lambda = \frac{\sigma^2}{\nu^2},
    \end{equation*}
    则 MAP 与带 $L_2$ 正则项的结构风险最小化完全一致, 对应解析解为
        \begin{equation*}
        \boldsymbol{w}^*_{MAP} = \arg\max_{\boldsymbol{w}} \log p(\boldsymbol{w}|\boldsymbol{X}, \boldsymbol{y}) = (\boldsymbol{X}\boldsymbol{X}^\top + \lambda \boldsymbol{I})^{-1}\boldsymbol{X}\boldsymbol{y}
        \end{equation*}
    \end{itemize}

  \item \textbf{估计思想总结}
    \begin{itemize}
    \item 线性回归中几种常见估计思想之间的关系可概括如下:
        \vspace{1em}
        \begin{center}
        \begin{tabular}{|c|c|c|}
        \hline
        & \textbf{无先验信息(唯数据论)} & \textbf{加入先验信息} \\
        \hline
        \textbf{函数测度视角(平方误差)} & 经验风险最小化(ERM) & 结构风险最小化(SRM) \\
        \hline
        \textbf{概率分布视角(概率统计)} & 最大似然估计(MLE) & 最大后验估计(MAP) \\
        \hline
        \multicolumn{3}{|c|}{最  优  参  数  的  闭  式  解: } \\
        \cline{1-3}
            $\boldsymbol{w}^*_{MLE / ERM}$ & \multicolumn{2}{c|}{$(\boldsymbol{X}\boldsymbol{X}^\top)^{-1}\boldsymbol{X}\boldsymbol{y}$} \\
        \cline{1-3}
            $\boldsymbol{w}^*_{MAP / SRM}$ & \multicolumn{2}{c|}{$(\boldsymbol{X}\boldsymbol{X}^\top + \lambda \boldsymbol{I})^{-1}\boldsymbol{X}\boldsymbol{y}$} \\
        \hline
        \end{tabular}
        \end{center}
    \end{itemize}
\end{enumerate}

(注: 详细定义和更多的内容请参考课程教材第二章、附录 D. )


\end{document}